\section*{الفصل \ref{ch:b1:comparative-planetology} --- علم الكواكب المقارن}

\begin{solution}{exo:b1:comparative-planetology:1}
$V = \tfrac43\pi(\qty{6.9911e7}{m})^3 = \qty{1.4313e24}{m^3}$ و
$\bar\rho = \num{1898.13e24}/\num{1.4313e24} = \qty{1326}{kg.m^{-3}}$، أي 1.33 مرة من كثافة الماء فقط، مع أن كتلة المشتري تفوق كتلة الأرض بنحو 318 مرة: فهو مكوّن في معظمه من أخفّ العناصر.
\end{solution}

\begin{solution}{exo:b1:comparative-planetology:2}
عطارد (نصف قطر \qty{2440}{km}، \qty{5429}{kg.m^{-3}}) والزهرة (\qty{6052}{km}، \qty{5243}{kg.m^{-3}}) صغيران وأكثف من الصخر: كوكبان صخريان، من صخر وحديد. وزحل (\qty{60268}{km}، \qty{687}{kg.m^{-3}}) ضخم وأخفّ من الماء: عملاق غازي، معظمه هيدروجين وهيليوم. وأورانوس (\qty{25559}{km}، \qty{1270}{kg.m^{-3}}) عملاق أصغر بكثير من زحل وأكثف منه بنحو الضعف: عملاق جليدي، ماؤه وميثانه وأمونياكه أثقل من الهيدروجين.
\end{solution}

\begin{solution}{exo:b1:comparative-planetology:3}
$T_\mathrm{eq} = (2601.3\times0.23/(4\times\num{5.6704e-8}))^{1/4} =
(\num{2.638e9})^{1/4} = \qty{226.6}{K}$. تتلقى الزهرة ضوءًا من الشمس أكبر بنحو 1.9 مرة من الأرض، لكن غيومها تعكس \qty{77}{\percent} منه: فهي لا تمتص إلا $2601.3\times0.23 = \qty{598}{W.m^{-2}}$، في مقابل $1361\times0.706 = \qty{961}{W.m^{-2}}$ للأرض.
\end{solution}

\begin{solution}{exo:b1:comparative-planetology:4}
من الأكثر انتظامًا إلى الأكثر تركّزًا: القمر، والمريخ، والأرض، والمشتري، وزحل. وكتلة نقطية في المركز لها $I = 0$، إذن $C = 0$.
\end{solution}

\begin{solution}{exo:b1:comparative-planetology:5}
$T = 15.945421\times86400 = \qty{1377684}{s}$;
$GM = 4\pi^2(\num{1.22187e9})^3/\num{1377684}^2 = \qty{3.794e16}{m^3.s^{-2}}$;
$M = \num{3.794e16}/\num{6.6743e-11} = \qty{5.685e26}{kg}$، أي في حدود
\qty{0.03}{\percent} من \qty{568.32e24}{kg}.
\end{solution}

\begin{solution}{exo:b1:comparative-planetology:6}
$S = 1361\times(149.6/228.0)^2 = \qty{585.9}{W.m^{-2}}$ عند المريخ، و
$1361\times(149.6/778.5)^2 = \qty{50.26}{W.m^{-2}}$ عند المشتري، وهو يتفق مع الجداول (والفرق الصغير عند المريخ يأتي من تقريب المسافات المتوسطة).
\end{solution}

\begin{solution}{exo:b1:comparative-planetology:7}
$S \propto d^{-2}$ و$T_\mathrm{eq} \propto S^{1/4}$، إذن $T_\mathrm{eq}
\propto d^{-1/2}$. وبالنسبة إلى المشتري، $T_\mathrm{eq} = (50.26\times0.657/(4\times
\num{5.6704e-8}))^{1/4} = \qty{109.8}{K}$. فالكوكب الذي يشعّ أكثر مما يمتص يجب أن يكون له مصدر حرارة داخلي: فالمشتري يطلق حرارة خزّنها منذ تكوّنه، أو يولّدها وهو ينكمش ببطء.
\end{solution}

\begin{solution}{exo:b1:comparative-planetology:8}
الزهرة: $737 - 226.6 = \qty{510}{K}$، أي \qty{225}{\percent} من $T_\mathrm{eq}$.
الأرض: $T_\mathrm{eq} = (1361\times0.706/(4\sigma))^{1/4} = \qty{255.1}{K}$،
$288 - 255.1 = \qty{33}{K}$، أي \qty{13}{\percent}. المريخ: $214 - 209.8 =
\qty{4}{K}$، أي \qty{2}{\percent}.
\end{solution}

\begin{solution}{exo:b1:comparative-planetology:9}
بوضع $\mathrm dm = 4\pi\rho r^2\,\mathrm dr$: $I = \int_0^R\tfrac23r^2\,4\pi
\rho r^2\,\mathrm dr = \tfrac{8\pi}{15}\rho R^5$، و$MR^2 = \tfrac43\pi\rho
R^5$، إذن $C = \tfrac{8}{15}\cdot\tfrac34 = \tfrac25$. وقشرة مجوّفة تتركّز كل كتلتها عند $r = R$: $I = \tfrac23MR^2$، و$C = \tfrac23$.
\end{solution}

\begin{solution}{exo:b1:comparative-planetology:10}
$\tfrac52C\bar\rho = 2.5\times0.35\times5429 = 4750.4$؛
$x^2 = (4750.4 - 3300)/(5429 - 3300) = 1450.4/2129 = 0.6812$، $x = 0.8254$،
نصف القطر اللبّي $0.8254\times2439.7 = \qty{2014}{km}$. $x^3 = 0.5623$ و
$\rho_c = 3300 + 2129/0.5623 = \qty{7086}{kg.m^{-3}}$. ويحوي اللبّ
$0.5623\times7086/5429 = \qty{73}{\percent}$ من الكتلة. ويبلغ لبّ عطارد \qty{83}{\percent} من نصف قطره، في مقابل $3485/6371 =
\qty{55}{\percent}$ للأرض: فعطارد في معظمه لبّ، مع وشاح صخري رقيق.
\end{solution}

\begin{solution}{exo:b1:comparative-planetology:11}
(أ) الدالة $t\mapsto t^4$ محدّبة، إذن
$\bigl(\tfrac{T_1 + T_2}{2}\bigr)^4 \le \tfrac{T_1^4 + T_2^4}{2} =
T_\mathrm{eq}^4$، ولا يتساويان إلا إذا كان $T_1 = T_2$. (ب) $T_1^4 =
2\times270^4 - 100^4 = \num{1.0529e10}$، إذن $T_1 = \qty{320}{K}$، ويكون المتوسط
$(320 + 100)/2 = \qty{210}{K}$، أي أدنى من $T_\mathrm{eq}$ بمقدار \qty{60}{K}.
\end{solution}

\begin{solution}{exo:b1:comparative-planetology:12}
$m = 4\pi R^2p/g = 4\pi\times(\num{6.0518e6})^2\times\num{9.2e6}/8.9 =
\num{4.602e14}\times\num{1.034e6} = \qty{4.76e20}{kg}$: أي 93 مرة كتلة
غلاف الأرض الجوي، و\num{9.8e-5} من كتلة الزهرة. وتحوي الأرض كمية
مماثلة من الكربون، ليس في هوائها، بل في الصخور الكربوناتية
(الحجر الجيري) وفي المحيط، حيث أتاح الماء السائل تخزينه.
\end{solution}

\begin{solution}{pb:b1:comparative-planetology:1}
\textbf{1.} $T = 0.31891\times86400 = \qty{27553.8}{s}$.

\textbf{2.} $GM = 4\pi^2\times(\num{9.378e6})^3/\num{27553.8}^2 =
\qty{4.289e13}{m^3.s^{-2}}$.

\textbf{3.} $M = \num{4.289e13}/\num{6.6743e-11} = \qty{6.426e23}{kg}$،
أي في حدود \qty{0.14}{\percent} من الكتلة الجدولية \qty{6.4169e23}{kg}.

\textbf{4.} $\num{1.06e16}/\num{6.426e23} = \num{1.6e-8}$: كتلة فوبوس مهملة.

\textbf{5.} $v = 2\pi a/T = 2\pi\times\num{9.378e6}/27553.8 =
\qty{2138}{m.s^{-1}}$، أي نحو \qty{2.1}{km.s^{-1}}.

\textbf{6.} $T = \qty{7.65}{h}$، أي أقصر بمقدار 3.2 مرة من دوران المريخ: فيدور فوبوس أسرع من دوران الأرض تحته، فيسبق الدوران. ويُرى من المريخ أنه يطلع من الغرب ويغرب في الشرق.

\textbf{7.} $V = \tfrac43\pi(\num{3.3895e6})^3 = \qty{1.6312e20}{m^3}$.

\textbf{8.} $\bar\rho = \num{6.426e23}/\num{1.6312e20} = 3939$، أي
\qty{3940}{kg.m^{-3}}.

\textbf{9.} $3940/5514 = 0.715$. فالأرض، الأكبر والأثقل، تضغط باطنها بقوة أكبر؛ وقلبها الحديدي حصة أكبر من الكوكب.

\textbf{10.} $g = GM/R^2 = \num{4.289e13}/(\num{3.3895e6})^2 =
\qty{3.73}{m.s^{-2}}$، أي \qty{38}{\percent} من جاذبية الأرض.

\textbf{11.} $v_\mathrm{esc} = \sqrt{2GM/R} = \sqrt{2\times\num{4.289e13}/
\num{3.3895e6}} = \qty{5.03}{km.s^{-1}}$.

\textbf{12.} $T_\mathrm{eq} = (586.2\times0.75/(4\times\num{5.6704e-8}))^{1/4}
= \qty{209.8}{K}$.

\textbf{13.} $214 - 209.8 = \qty{4}{K}$: أثر دفيئ ضعيف.

\textbf{14.} $m = p/g = 636/3.733 = \qty{170}{kg.m^{-2}}$؛ أما للأرض
$101400/9.82 = \qty{10326}{kg.m^{-2}}$، أي أكثر بنحو 61 مرة.

\textbf{15.} $4\pi R^2\times170.4 = \num{1.444e14}\times170.4 =
\qty{2.46e16}{kg}$.

\textbf{16.} فوق كل متر مربع من المريخ غاز أقل ستين مرة مما على الأرض، لذلك يكون امتصاص تحت الحمراء ضعيفًا، مع أنه بثاني أكسيد الكربون؛ وهواء المريخ يحوي بخار ماء قليلًا جدًا، وهو الماص الرئيسي الآخر.

\textbf{17.} لو كان المريخ منتظمًا لكان $C = 0.4$؛ والقيمة المقيسة 0.366 تدلّ على أن كثافته تتزايد نحو المركز: للمريخ لبّ كثيف.

\textbf{18.} $\bar\rho = \rho_m + (\rho_c - \rho_m)x^3$.

\textbf{19.} $C = \tfrac25[\rho_m + (\rho_c - \rho_m)x^5]/\bar\rho$.

\textbf{20.} من السؤالين 18 و19، $(\rho_c - \rho_m)x^3 = \bar\rho - \rho_m$ و
$(\rho_c - \rho_m)x^5 = \tfrac52C\bar\rho - \rho_m$؛ وبالقسمة نحصل على $x^2$.

\textbf{21.} $\tfrac52C\bar\rho = 2.5\times0.366\times3940 = 3605.1$;
$x^2 = (3605.1 - 3500)/(3940 - 3500) = 105.1/440 = 0.2389$, $x = 0.4887$.
$x^3 = 0.1167$ و$\rho_c = 3500 + 440/0.1167 = \qty{7269}{kg.m^{-3}}$.

\textbf{22.} $x^3\rho_c/\bar\rho = 0.1167\times7269/3940 = 0.215$، أي
\qty{21.5}{\percent} من الكتلة.

\textbf{23.} مع $\rho_m = 3450$: $x^2 = 155.1/490 = 0.3165$، $x = 0.5626$،
ونصف قطر اللبّ \qty{1907}{km}، و$\rho_c = \qty{6202}{kg.m^{-3}}$. ومع $\rho_m =
3550$: $x^2 = 55.1/390 = 0.1413$، $x = 0.3759$، و\qty{1274}{km}،
$\rho_c = \qty{10894}{kg.m^{-3}}$. فتغيّر قدره \qty{1.4}{\percent} في كثافة الوشاح المفروضة يحرّك نصف قطر اللبّ بعدة مئات من الكيلومترات: فالكثافة المتوسطة و$C$ وحدهما لا يحدّدان لبّ المريخ بدقة، ولهذا أُرسلت مقاييس الزلازل إليه.

\textbf{24.} $0.4887\times3389.5 = \qty{1657}{km}$: نصف قطر لبّ المريخ من النموذج هو \textbf{نحو \qty[text-series-to-math=true]{1660}{km}}، ضمن المدى الزلزالي الحديث، \qty{1650}{km} إلى \qty{1675}{km}، ودون \qty{1830}{km} في الجدول الكوكبي بكثير ---
وهو توافق يبيّن السؤال 23 أنه جزئيًا محض حظّ في اختيار $\rho_m$.
\end{solution}
